\section{Scope: the unresolved class and the theorem proved here}
The elementary prime reduction leaves $p\equiv1\pmod{24}$. The preceding
receiving-determinant theorem applies to every existing integral witness at
$p\equiv1\pmod{12}$, so already includes this class; it does not assert that a
witness exists. We now return to the original arithmetic source. The main
result is the following source-changing theorem, rather than an improvement
of the determinant constant.

\begin{theorem}[Mixed prime-shell trace return]\label{thm:mixed}
Let $p\equiv1\pmod{24}$ be prime. Suppose $a=cq$ where $c\mid6$, $q>3$ is prime,
and $p/4<a<p/2$. If $y,z\in\mathbb Q_{>0}$ satisfy
\[
 \frac4p=\frac1a+\frac1y+\frac1z,\qquad y+z\in\mathbb Z,
\]
then the formulas in Sections~\ref{sec:auto}--\ref{sec:mixed} construct an
original integral middle state at the same $p$. All input orderings and
original divisor exponents have explicit inverses or stated fibres. If the
input tails are nonintegral, the returned middle state's smaller residual or
cofactor is strictly smaller than the input residual. In the exterior case
there is additionally an original exterior return with smaller first denominator.
\end{theorem}

The theorem concerns positive denominators in the complete E/M atlas; no extra
Koide-chamber condition is imposed on the retained ray. The theorem does not
assert the existence of its initial rational trace source.
At $p=2521$ and $3361$, the entire displayed family of initial shells is empty.
We also give an infinite obstruction to extending the fixed-ray return beyond
its precise nine-ray domain. The reductions are unconditional on their stated
source sets; none supplies a universal occupancy hypothesis.

The original exterior congruences are Yamamoto's Lemma~2, equation~(4), p.~38
\cite{Yamamoto}. The Type I/II context is Elsholtz--Tao, Section~2, Propositions
2.2 and 2.6 \cite{ET}. Their identities are rederived below. The reciprocal
trace source and the distinct fixed-word deletion are inherited from the
preserved workbench tranches \cite{Trace,Word}. No historical-priority conclusion
is made for the present continuation.

\section{The full rational trace source and original coordinates}\label{sec:source}
Fix a prime $p\equiv1\pmod4$, $p/4<a<p/2$, and $R=4a-p$. Thus
$3\le R<p$, $R\equiv3\pmod4$, and $\gcd(pa,R)=1$.
For every original divisor $u\mid a^2$, put
\[
 g=\gcd(a,u),\quad (h,r,s)=(g^2/u,u/g,a/g),\quad
 a=hrs,\quad u=hr^2,\quad \gcd(r,s)=1.
\]
If $a=\prod\ell^{e_\ell}$ and $u=\prod\ell^{f_\ell}$, the complete box is
$0\le f_\ell\le2e_\ell$; its centred coordinate is
$\beta_\ell=f_\ell-e_\ell=v_\ell(r)-v_\ell(s)$.
Before imposing integrality of the tails, the two positive rational returns are
\begin{align}
 E:\quad &(a,hs\kappa,phr\kappa),&\kappa&=(pr+s)/R,\label{eq:E}\\
 M:\quad &(a,phs\lambda,phr\lambda),&\lambda&=(r+s)/R.\label{eq:M}
\end{align}
Equivalently, they are
\[
 E:\left(a,\frac{pa+a^2/u}{R},\frac{pa+p^2u}{R}\right),\qquad
 M:\left(a,\frac{p(a+a^2/u)}R,\frac{p(a+u)}R\right).
\]
Direct common-denominator multiplication proves the reciprocal identities.
The full gates and the ordered divisor inverses are
\begin{align*}
 E:&\ R\mid4u+1\ \Longleftrightarrow\ R\mid pr+s,
 &u&=a^2/(Ry-pa),\\
 M:&\ R\mid p+4u\ \Longleftrightarrow\ R\mid r+s,
 &u&=pa^2/(Ry-pa).
\end{align*}
All cancelled factors are units modulo $R$.

\begin{lemma}[Exact trace relaxation]\label{lem:trace}
Let $G_E=4u+1$ and $G_M=p+4u$. The rational tails have integral sum precisely
when $R\mid G_c^2$. Their common reduced denominator is
\[
 d=R/\gcd(R,G_c),\qquad d^2\mid R.
\]
For E, the trace condition is also $R\mid(pr+s)^2$, and
$\gcd(R,G_E)=\gcd(R,pr+s)$. For M,
\(\gcd(R,G_M)=\gcd(R,r+s)\). A proper trace source means $d>1$.
Every positive rational pair at the prescribed integral $a$ with integral
sum occurs in these coordinates, up to the separately retained exterior swap.
\end{lemma}
\begin{proof}
The two tail sums are $(a+pu)^2/(Ru)$ and $p(a+u)^2/(Ru)$.
Their numerators and denominators, using $u\mid a^2$, give the stated square
conditions and common denominator. Alternatively, for arbitrary rational tails
with integral sum $N$, their reduced denominators coincide at $d$, since their
sum is integral. Their product is $Npa/R$. Prime-by-prime valuation shows
\[
 2v_\ell(d)=v_\ell(R)-v_\ell(N)\quad\text{when }\ell\mid d,
\]
while \(v_\ell(N)\ge v_\ell(R)\) when \(\ell\nmid d\). Hence
$d^2=R/\gcd(R,N)$. Thus $d\mid R$, and
\[
 b=Ry-pa,\qquad c=Rz-pa
\]
are positive integers with $bc=(pa)^2$. Their $p$-valuation pair is $(0,2)$,
$(1,1)$ or $(2,0)$. For $(0,2)$ set $u=a^2/b$; for $(2,0)$ first swap the
tails and retain that bit; for $(1,1)$ set $u=pa^2/b$. The resulting integer
$u$ divides $a^2$ and gives exactly \eqref{eq:E} or \eqref{eq:M}.
Finally,
\[
 pr+s\equiv s(4u+1)\pmod R,\qquad
 p(r+s)\equiv s(p+4u)\pmod R
\]
proves both gcd statements.
\end{proof}

\begin{lemma}[Original characters, including proper traces]\label{lem:char}
On every trace source,
\[
 E:\left(\frac up\right)=-1,\qquad
 M:\left(\frac up\right)=-\left(\frac ap\right).
\]
\end{lemma}
\begin{proof}
For odd $\ell\mid a$, quadratic reciprocity and $R\equiv-p\pmod\ell$ give
$(\ell/R)=(p/\ell)=(\ell/p)$. If $2\mid a$, the supplementary law gives the
same equality from $R\equiv-p\pmod8$. Multiplication through the original
exponents gives $(u/R)=(u/p)$. At each prime \(\ell\) dividing \(R\), the
square gate implies only the primewise statement \(\ell\mid G_c\), or
equivalently \(\operatorname{rad}(R)\mid G_c\); it does not imply
\(R\mid G_c\). This primewise statement is exactly what the Jacobi product
uses. Thus the Jacobi value for E is $(-1/R)=-1$;
for M it is $(-p/R)=-(R/p)=-(a/p)$. Prime multiplicities in $R$ are retained
in these Jacobi products.
\end{proof}

\section{A ray-preserving return with its full integer coefficient}\label{sec:ray}
For an E trace source define $m=4rs$, $T=pr+s$ and
$d=R/\gcd(R,T)$. Retaining the primitive ray means retaining both $r$ and $s$,
not retaining the old numerical divisor $u$.

\begin{theorem}[Complete divisor-deletion criterion]\label{thm:ray}
For \(k\in\mathbb Z_{>0}\) with $k\mid R$, a ray-preserving integral exterior
return at $R'=R/k$ exists
exactly when
\begin{equation}\label{eq:delete}
 d\mid k,\qquad k\equiv1\pmod m.
\end{equation}
The literal target is
\[
 h'=(p+R/k)/m,\quad a'=h'rs,\quad u'=h'r^2,\quad
 \kappa'=T/(R/k)=k\kappa.
\]
For a proper source, $a'<a$. At every prime $\ell$ the centred exponent is
unchanged: $v_\ell(u')-v_\ell(a')=v_\ell(r)-v_\ell(s)$.
\end{theorem}
\begin{proof}
The full E gate is $R/k\mid T$, exactly $d\mid k$. Since
\(\gcd(R,a)=1\), \(R\) is odd and \(rs\mid a\), one has
\(\gcd(R,m)=1\); hence every \(k\mid R\) is a unit modulo \(m\).
The integer shell condition is therefore equivalent to
\[
 m\mid p+R/k
 \Longleftrightarrow m\mid pk+R
 \Longleftrightarrow m\mid p(k-1)
 \Longleftrightarrow k\equiv1\pmod m.
\]
The middle equivalence uses \(m\mid p+R\), and the last uses
\(\gcd(p,m)=1\), which follows from \(rs\le a<p\). It also gives
$R'\equiv3\pmod4$,
so $R'\ge3$ and the target is in the original first-half range. All target
coordinates and their ordered inverse follow from \eqref{eq:E}: explicitly,
\[
 \frac{(a')^2}{u'}=h's^2,\qquad
 \gcd(a',u')=h'r,\qquad
 R'y'-pa'=h's^2,\qquad
 u'=\frac{(a')^2}{R'y'-pa'}.
\]
Properness
forces $k>1$. The centred equality follows from the two factorizations of $a,u$.
\end{proof}

The complete number of deletion words is the integer coefficient
\begin{equation}\label{eq:coefficient}
 [d^{-1}]\prod_{\ell^e\parallel R/d}
 \bigl([1]+[\ell]+\cdots+[\ell^e]\bigr)
 \quad\hbox{in }\mathbb Z[(\mathbb Z/m\mathbb Z)^\times].
\end{equation}
Each exponent vector occurs once; no binomial coefficients are substituted
for its original multiplicity. The inverse reads the retained old $R$ and $k$.
A fixed-word return instead uses modulus $4K(u)$, where
$K(u)=\prod\ell^{\lceil v_\ell(u)/2\rceil}$ \cite{Word}. The two conditions
are different: the ray return changes $u$ while preserving the centred exponents.

\section{Nine automatic rays and their complete middle images}\label{sec:auto}
The unit squares modulo $m=4rs$ are all one if and only if $m\mid24$.
On any odd prime-power factor other than $3$ to exponent one, the unit $2$
has $2^2-1=3$ nonzero in that factor. On a factor $2^e$ with $e\ge4$, the
unit $3$ has $3^2-1=8$ nonzero. CRT extends either witness to a unit modulo
the full $m$. Conversely every odd integer has square one modulo eight,
and every unit has square one modulo three. CRT gives the converse. Equivalently, $rs\mid6$.
The nine primitive ordered rays are
\[
 (1,1),(1,2),(1,3),(1,6),(2,1),(2,3),(3,1),(3,2),(6,1).
\]
For every such E trace source, write $R=\delta t^2$, where $\delta$ is the
positive squarefree part: prime exponents are reduced modulo two. This is not
the radical, and $\gcd(\delta,t)$ need not be one. Since $R\mid T^2$,
$\delta\mid T$, and $t^2\equiv1\pmod m$. Therefore $k=t^2$ satisfies
\eqref{eq:delete}, giving a full E state at $R'=\delta$.

When $p\equiv1\pmod{24}$, one has $p\equiv1\pmod m$ and
$\delta\equiv-1\pmod m$. Set
\begin{equation}\label{eq:Mreturn}
 H=(\delta+1)/m,\qquad J=T/\delta,\qquad
 (h_M,r_M,s_M,\lambda_M)=(H,s,J,r).
\end{equation}
Then
\[
 a_M=HsJ,\quad u_M=Hs^2,\quad R_M=(s+J)/r,\quad Q_M=4Hsr-1=\delta.
\]
These are integers: $\delta J=pr+s$ and $\delta\equiv-1\pmod r$ imply
$r\mid s+J$. A common factor of $s,J$ would divide $p$ or $r$, which is
impossible. Directly,
\[
 4a_M-R_M=((\delta+1)J-s-J)/r=p.
\]
Furthermore
\[
 0<R_M\le p/3+4s/(3r)<p,
\]
because $p>2hrs\ge2s/r$. Thus the original ordered middle denominators are
\begin{equation}\label{eq:MDenom}
 (HsJ,\ pHJr,\ pHsr).
\end{equation}
If the source is proper, $t>1$ and
$\min(R_M,Q_M)\le\delta<R$. This is a cofactor decrease, not necessarily
a decrease in $R_M$ or $a_M$. The accompanying E return always has $a'<a$.

\subsection{Exact squarefree fibres}
For fixed $p$ and an automatic ray, let $A$ be the odd part of $T=pr+s$.
The complete target index set consists of squarefree divisors
\[
 \delta\mid\operatorname{rad}(A),\qquad \delta\equiv-1\pmod m.
\]
Here the symbol $\operatorname{rad}(A)$ only specifies the squarefree divisor
set; it does not replace the source's squarefree part. Every such $\delta<p$:
$J\equiv-(r+s)\pmod m$ gives
$J\ge m-r-s\ge r+s$, and $(pr+s)/(r+s)<p$.
The full trace fibre is
\begin{equation}\label{eq:fibre}
 R=\delta t^2,
 \qquad t\mid A/\delta,
 \qquad \delta t^2<p.
\end{equation}
To prove this, compare exponents in $R\mid T^2$. At a prime of $\delta$, the
exponent of $t$ is at most $v_\ell(A)-1$; elsewhere it is at most $v_\ell(A)$.
These conditions are also sufficient. Here \(\gcd(A,m)=1\): for every prime
\(\ell\mid r\), \(T\equiv s\not\equiv0\pmod\ell\), while for
\(\ell\mid s\), \(T\equiv pr\not\equiv0\pmod\ell\). Thus every allowed \(t\)
is a unit modulo \(m\), and its square is one. The congruence of $R$ is
therefore preserved by all such squares modulo $m$. The full original hits within this fibre are
exactly those additionally satisfying $R\mid T$.

Retaining $t$ in \eqref{eq:fibre} gives an inverse. From the marked middle
image, recover $r=\lambda_M$, $s=r_M$, $J=s_M$ and
$\delta=4h_Mr_M\lambda_M-1$, then reconstruct the old residual and $h$.
Let \(e_{\delta,t}\) denote a source basis vector and \(f_\delta\) the
corresponding target basis vector. Forgetting \(t\) gives the split map
\[
 \pi:\mathbb Z[\{(\delta,t)\}]\longrightarrow\mathbb Z[\{\delta\}],
 \qquad e_{\delta,t}\longmapsto f_\delta,
 \qquad \sigma(f_\delta)=e_{\delta,1}.
\]
Its kernel has basis \(e_{\delta,t}-e_{\delta,1}\) for \(t>1\) in the
original bounded fibre.
Neither surjectivity nor the section asserts that the target index set is
nonempty. At $p=1129$, ray $(3,1)$ and $\delta=11$ have the two words $t=1,7$;
its integer coefficient two is not absent merely because it reduces to zero
modulo two.

\section{Proof of the mixed prime-shell theorem}\label{sec:mixed}
By Lemma~\ref{lem:trace}, normalize the given rational pair to an original
E or M trace word $u\mid a^2$, retaining its input orientation. Both two and
three are quadratic residues at $p\equiv1\pmod{24}$. If $q$ were a residue,
every divisor of $a^2$ and $a$ itself would have positive character, contradicting
Lemma~\ref{lem:char}. Thus $(q/p)=-1$.

For E, the exponent of $q$ in $u$ is odd, hence exactly one. Write
$u=qw$, $w\mid c^2$. The gcd normalization has $rs\mid c\mid6$:
for each prime of the squarefree integer $c$, its exponent zero or two in $w$
contributes that prime to exactly one of $r,s$, and exponent one contributes
 it to neither. The automatic ray construction therefore gives
\eqref{eq:Mreturn}--\eqref{eq:MDenom}. If the arbitrary input tails required
the exterior swap in Lemma~\ref{lem:trace}, replace the canonical returned
middle word \(u_M=Hs^2\) by its new-shell complement
\[
 \frac{a_M^2}{u_M}=HJ^2;
\]
this interchanges the two middle tails and restores the retained input
orientation.

For M, Lemma~\ref{lem:char} forces an even exponent of $q$. If it is two,
replace $u$ by its original complement $a^2/u$ and retain that bit; otherwise
leave it unchanged. The resulting $u_0$ divides $c^2$, hence divides $36$.
This complement preserves the square gate because
\[
 u\left(p+4a^2/u\right)-a(p+4u)=(a-u)R,
\]
and \(a,u\) are units modulo \(R\). Thus $4K(u_0)\mid24$. With
$R=\delta t^2$ one has
\[
 \delta\mid p+4u_0,\qquad
 t^2\equiv1\pmod{4K(u_0)},\qquad
 p+\delta\equiv0\pmod{4K(u_0)}.
\]
Put $a'=(p+\delta)/4$. These equations prove $u_0\mid a'^2$ and the full M
gate at the original residual $\delta$. Recompute the gcd normalization and
use \eqref{eq:M}. If the old complement bit was one, use $u'=a'^2/u_0$ to
restore the corresponding tail ordering. The new exponent box, not the old
factor inventory, is used in this reversal. Marking the bit
\(\varepsilon\in\{0,1\}\), the inverse reads
\[
 u_0=\begin{cases}u_{\rm out},&\varepsilon=0,\\
 a'^2/u_{\rm out},&\varepsilon=1,\end{cases}
 \qquad
 u=\begin{cases}u_0,&\varepsilon=0,\\
 a^2/u_0,&\varepsilon=1,\end{cases}
\]
with the retained square factor \(t\) recovering \(R=\delta t^2\) and
\(a=(p+R)/4\). A proper input has $\delta<R$,
so $a'<a$. This proves Theorem~\ref{thm:mixed}.

The excluded carrier values $q=2,3$ provide no trace source at these primes:
all available prime factors would be residues. Without the $p\equiv1\pmod{24}$
hypothesis the carrier conclusion can fail, as the exact $p=37,a=10,u=2$ E
state demonstrates.

\section{Examples, source fibres and indispensable negative controls}\label{sec:examples}
At the hard prime $p=1129$, let
\[
 a=417=3\cdot139,\quad R=539=11\cdot7^2,\quad u=1251,
 \quad(h,r,s)=(139,3,1).
\]
The positive rational E identity has tails $6116/7$ and $20714892/7$, with
integer sum $2960144$. The complete original shell has both gates empty.
Squarefree ray descent gives
\[
 (a',R',u',h',r',s',\kappa')=(285,11,855,95,3,1,308),
\]
with denominators $(285,29260,99103620)$. The middle return is
\[
 (a_M,R_M,u_M,h_M,r_M,s_M,\lambda_M)
 =(308,103,1,1,1,308,3),
\]
\[
 \frac4{1129}=\frac1{308}+\frac1{1043196}+\frac1{3387}.
\]
The unchanged word $1251$ does not divide $285^2$, and its entire previous
fixed-word deletion set is empty. This is a genuinely different transport.

At $p=2689$, $a=1138=2\cdot569$, $R=1863=23\cdot3^4$, $u=569$, both old
gates again fail. Here $\delta=23$, while the radical of $R$ is $69$ and
is $1\pmod4$: using it would destroy the original shell. At $p=6841$,
\[
 a=2454=6\cdot409,\quad R=2975=7\cdot17\cdot5^2,\quad u=1636
\]
has rational tails $28221/5$ and $128706574/5$. Its squarefree part is the
composite $119$, not a selected prime factor. The middle return is
\[
 (a_M,R_M,u_M,h_M,r_M,s_M,\lambda_M)=(1725,59,45,5,3,115,2),
\]
\[
 \frac4{6841}=\frac1{1725}+\frac1{7867150}+\frac1{205230}.
\]
Every original and returned prime exponent is serialized in the certificate.
The $p=1009,a=321,u=9$ M example returns $(255,686120,24216)$ and cross-checks
the earlier fixed-word theorem rather than supplying a new example of that theorem.

\subsection{The factor \texorpdfstring{$4q$}{4q} cannot be added to this deletion rule}
At the hard prime $p=9601$, the M source
\[
 a=2956=4\cdot739,\quad R=2223=3^2\cdot13\cdot19,\quad u=8
\]
has tails $14190278/3,38404/3$ and an integral sum. The squarefree target is
$\delta=247$, $a'=2462=2\cdot1231$. Its complete divisor-residue set is
\[
 \{1,2,4,16,32,64,231,239,243\}\pmod{247}.
\]
It contains neither the E target $185$ nor the M target $8$. This refutes
this automatic squarefree deletion at $4q$, not all possible source-changing
repairs at that shell. An original M state at $a=2405,R=19,u=65$ is separately
certified.

\subsection{The initial-source implication is still missing}
Exhaustion at $p=2521$ and $3361$ finds no rational integral-trace pair in any
of the shells $cq$ specified by Theorem~\ref{thm:mixed}. All those prime-shell
records and empty square gates are supplied. At $3361$, every one of the nine
automatic E rays also has no initial trace state. Their linear forms factor as
\[
\begin{array}{c|l@{\qquad}c|l}
(r,s)&pr+s&(r,s)&pr+s\\\hline
(1,1)&2\cdot41^2&(1,2)&3\cdot19\cdot59\\
(1,3)&2^2\cdot29^2&(1,6)&7\cdot13\cdot37\\
(2,1)&3^4\cdot83&(2,3)&5^2\cdot269\\
(3,1)&2^2\cdot2521&(3,2)&5\cdot2017\\
(6,1)&7\cdot43\cdot67&&
\end{array}
\]
None has a divisor in its required class $-1\pmod{4rs}$. The complete fibre
formula therefore excludes its trace sources too. This is not an ES
counterexample: $a=841=29^2$, $R=3$, $u=29$ is an original E state (and M state).
The carrier in the failed mixed form would be composite.

\section{Optimality of the automatic-ray guarantee}\label{sec:optimal}
\begin{theorem}[An infinite original obstruction outside the nine rays]
Fix coprime positive $r,s$ with $rs\nmid6$. There are infinitely many hard
primes $p\equiv1\pmod{840}$ with an original proper E trace source on this ray
for which every ray-preserving residual-divisor return is absent. The sources
may simultaneously have no fixed-word residual-divisor return, and their E
complements may fail even the relaxed trace gate. These primes have separately
constructed integer middle solutions. No global ES counterexample is claimed.
\end{theorem}
\begin{proof}
Put $m=4rs$. Since $m\nmid24$, choose a unit $b$ with $b^2\not\equiv1\pmod m$;
change its sign if necessary so that $b\equiv1\pmod4$. By Dirichlet's theorem
\cite[Theorem 18.1]{Dirichlet}, choose distinct primes
\[
 t\equiv b\pmod m,\qquad \delta\equiv-b^{-2}\pmod m,
\]
both exceeding $\max(7,r,s,|r^2-s^2|)$. Then $t\equiv1\pmod4$,
$\delta\equiv3\pmod4$. Set $R=\delta t^2$ and $L=\operatorname{lcm}(840,m)$.
Neither \(t\) nor \(\delta\) divides \(L\): both exceed seven and both are
units modulo \(m\). Choose a prime $e\equiv3\pmod4$ coprime to $LR$. The CRT conditions
\begin{equation}\label{eq:CRT}
 p\equiv1\pmod L,\qquad pr+s\equiv\delta t\pmod R,
 \qquad p\equiv-1\pmod e
\end{equation}
form a reduced class: the moduli are pairwise coprime, $r$ is a unit modulo
$R$, and \(\delta t-s\equiv-s\) is nonzero modulo both $\delta$ and $t$
because they exceed \(s\). Dirichlet supplies
infinitely many primes in it.

Choose them large enough that $p>R$ and $p>sR^2/(4r)$. Since $R\equiv-1\pmod m$,
$h=(p+R)/m$, $a=hrs$ and $u=hr^2$ give an original proper E trace source
with integral shell coordinates. Writing
\(pr+s=\delta t+n\delta t^2=\delta t(1+nt)\) shows exactly that
\(\gcd(R,pr+s)=\delta t\), so its trace defect is $d=t$. The possible deleting
factors in \eqref{eq:delete} are $t,t^2,\delta t,\delta t^2$. The last two are
$3\pmod4$; the first two are $b,b^2\not\equiv1\pmod m$. There is no ray return.

Moreover $u>R^2/16$, so $4K(u)\ge4\sqrt u>R$. The fixed-word budget follows
directly from \(K(u)\mid a\): both source and target shell conditions force
\(k\equiv1\pmod{4K(u)}\). No divisor \(k\mid R\) with \(d=t>1\) can satisfy
that congruence when \(4K(u)>R\). For the complement $u^*=hs^2$,
\[
 4u^*+1\equiv (r^2-s^2)r^{-2}\not\equiv0\pmod\delta.
\]
Here $r\ne s$, since coprimality and equality would give the automatic ray
$(1,1)$. Thus the complement is not even a trace source.
Finally $e\mid p+1$ gives $H=(e+1)/4$, $J=(p+1)/e$ and the original middle
state $(h_M,r_M,s_M,\lambda_M)=(H,1,J,1)$, with denominators
$(HJ,pHJ,pH)$. Its residual is \(R_M=J+1\), and
\[
 0<R_M\le\frac{p+1}{3}+1=\frac{p+4}{3}<p,
\]
so it is in the first-half range.
\end{proof}

This proves an exact classification of the \emph{universal repair rule on a
supplied trace ray}: it works for all such sources exactly when $rs\mid6$.
It does not say that nonautomatic rays never succeed. The finite progression
records retain reduced residue, modulus and lower threshold; their sample
integers are not described as prime values.

\subsection{The older third Farey ray has a completely failed trace example}
Take
\[
 p=6975049201,\quad r=5,\quad s=3,\quad
 R=378359=71\cdot73^2,
\]
\[
 h=116257126,\quad a=1743856890,\quad u=2906428150.
\]
This prime lies in the reduced progression
$p=3479012041+3496037160n$ at $n=1$. Its original proper E trace has tails
\[
 \frac{2346804477773328}{73},\qquad
 \frac{27281794495993456209184880}{73}.
\]
The complete linear-form factorization is
\[
 5p+3=34875246008=2^3\cdot71\cdot73\cdot841097.
\]
Every required residual is odd. The complete odd-divisor residue set is
\[
 \{1,7,11,13,17,23,31,41\}\pmod{60},
\]
which omits $59$. Thus the \emph{entire genuine ray} $(5,3)$ fails here, not
just the residual-divisor return. The 243 divisor words \(u\mid a^2\), for
$a=2\cdot3\cdot5\cdot293\cdot198391$, also give empty full E and M gates at
the specified old shell. These are exact finite exhaustions in the certificates.

Primality is checked by two methods: trial division through $\lfloor\sqrt p\rfloor$
and, separately, the complete-order certificate
\[
 p-1=2^4\cdot3\cdot5^2\cdot7\cdot830363,\qquad\hbox{base }11.
\]
Each listed factor is prime; the certificate checks $11^{p-1}=1\pmod p$ and
$\gcd(11^{(p-1)/q}-1,p)=1$ for every listed $q$. These force order $p-1$ modulo
any prime divisor of $p$, proving primality. The same prime has the independently
available middle state
\[
 (h_M,r_M,s_M,\lambda_M)=(3,1,634095382,1),
\]
\[
 (x,y,z)=(1902286146,\ 6975049201\cdot1902286146,\ 20925147603).
\]
No least-prime claim or full verification through this large isolated prime is made.

\section{Reproduction, scope and the remaining implication}
The main standard-library verifier enumerates every original first-half square
divisor at primes $p\equiv1\pmod{24}$ through the declared bound. The separate
checker imports no main or predecessor code and enumerates primitive $h,r,s$
parameters instead. Both test the square gates, original full gates, centred
exponent transport, middle returns, fibres, CRT records and explicit obstructions.
\begin{verbatim}
python verify.py --bound 10000 --out certificates
python -O verify.py --bound 10000 --out certificates_optimized
python check_independent.py --input certificates \
  --out certificates/independent.json
python repair.py --p 1129 --a 417 --y 6116/7 --z 20714892/7
\end{verbatim}
The complete run contains 143 primes, 171732 shells, 5616298 original divisor
words, 5788 E traces and 5678 oriented M traces. Of these, 3637 E and 3910 M
are integral. The mixed family has 806 trace states, 254 proper; the automatic
rays have 1230 E traces, 409 proper. These counts overlap as specified source
subsets and are not added as new prime coverage. Thirty-four nonautomatic-ray
progression templates with $r,s\le8$ are arithmetic CRT certificates, not lists
of independently certified prime sample values.

The finite tests corroborate the implementation. The general theorems have
the preceding coordinate proofs and the explicitly reduced Dirichlet input.
No Lean build or independent human mathematical review is claimed. The current
result completes the stated trace-to-integer return; it does not prove
\[
 \forall p\equiv1\pmod{24}\text{ prime},\quad
 \exists\text{ an occupied original E or M coefficient}.
\]
In particular the split coefficient map and its integral section in
\eqref{eq:fibre} do not establish a first source. The negative controls specify
why both a supplied-ray assumption and a supplied mixed prime-shell trace
assumption remain substantive. The programme therefore remains at Turn~7.
